\section*{Part II: How the Mind Makes Knowledge Possible}

\subsection*{13. Prior Knowledge Is the Medium into Which New Ideas Arrive}

No explanation enters an empty mind. An engineer hears ``retries with backoff'' and connects it to incidents, latency budgets, and network failures. A novice sees a confusing phrase. A mathematician sees the structure of a proof because thousands of earlier examples have made the relevant abstractions available.

Prior knowledge determines what the learner notices, what seems plausible, and how many interacting details can be handled. It may be accurate, incomplete, fragmented, or actively misleading. A person who believes a high test accuracy always means a good model will assimilate a sophisticated lesson into the wrong mental model unless the misconception is confronted.

The practical requirement is \textbf{diagnosis before major instruction}. Not a prestige-based label or a self-rating, but a small authentic task: explain this output, diagnose this error, predict this result. The response reveals which prerequisites are present and which conceptual structures need repair [R1].

\subsection*{14. Expertise Reorganizes the Problem the Novice Sees}

An expert is not merely a novice with more facts. Familiar patterns become meaningful units. A seasoned developer sees a design with 3 failure domains; a novice sees dozens of disconnected lines of code. A mathematician recognizes that a proof requires an invariant; a novice searches through formulas at random.

This pattern recognition can make expert reasoning fast, but it also creates \textbf{the expert blind spot}. An instructor may skip the steps that their own knowledge has compressed. ``Obviously, choose the appropriate validation split'' hides several difficult judgments about how the data arise and how the model will be used.

To teach effectively, the expert must unpack their own perception. Ask them to annotate not only what they decided, but which cues they noticed, which alternatives they rejected, and which assumptions would change the answer. The teaching value often lies in exposing distinctions that have become invisible to the instructor.

\subsection*{15. Attention Is the Gateway, but Not a Complete Explanation}

Attention determines which information receives enough processing to become usable. Learners divide attention among the instructor, code, notifications, anxieties, competing tabs, and their own predictions. An elegant video can capture attention without building knowledge. A boring-looking worksheet can demand deep attention to an important relationship.

The founder must design \textbf{attention toward the right cognitive object}. Signaling and reducing irrelevant detail can help. But asking whether learners ``paid attention'' is inadequate. The decisive question is: \emph{To what did their attention lead?} Did they merely track the teacher's explanation, or did they independently predict, construct, compare, and revise?

Learning environments should protect periods of focused work while respecting accessibility and varied circumstances. An adult learner returning from a demanding shift may need a manageable segment and clear next action, not a theatrical demand for uninterrupted enthusiasm.

\subsection*{16. Working Memory Is Constrained; Long-Term Knowledge Changes Effective Capacity}

Unfamiliar tasks require active manipulation of elements in working memory, which is limited. Experts can appear to juggle more information because relevant relationships are already organized in long-term knowledge. The novice must consciously process every line; the expert recognizes a higher-order pattern.

Cognitive load theory distinguishes the demands inherent in a task from avoidable demands imposed by poor presentation. Its practical value is not the promise of a simple numerical ``load meter.'' It is a warning against making a learner solve the content problem \emph{and} the communication problem at once. Worked examples, pretraining, reduced visual search, and thoughtfully segmented explanations can help novices [R2].

The opposite mistake is to strip all difficulty away. Learning to design systems eventually requires handling interacting constraints. The right progression reduces irrelevant complexity early, then deliberately reintroduces the complexity that belongs to the real professional task.

\subsection*{17. The Expertise-Reversal Effect: Guidance Has a Half-Life}

A detailed worked example can help a beginner who lacks the schema to solve a problem. The same example may irritate or slow a knowledgeable learner who could solve the task faster independently. Research on the \textbf{expertise-reversal effect} shows that instructional supports do not have fixed value independent of the learner's knowledge [R3].

This suggests a general rule: \textbf{diagnose readiness, scaffold what is missing, and fade the scaffold when it becomes redundant}. It is not sufficient to separate people permanently into ``beginner'' and ``advanced'' tracks; an experienced backend engineer may be advanced in software architecture and novice in statistical experimental design.

The founder should therefore avoid selling personalization as a mysterious algorithm. Often the most valuable adaptation is relatively simple: reveal a weak prerequisite, provide a strong example, ask for an attempt, then change the level of guidance in response.

\subsection*{18. Models, Schemas, and Abstraction Are the Foundations of Useful Understanding}

A schema is an organized representation that helps someone interpret new information. It is not merely a collection of associated facts. A useful schema can generate predictions and support decisions.

Consider conditional probability. A learner might remember Bayes' formula but misunderstand the base rate. Another can reason with natural frequencies, draw a contingency table, and recognize why a seemingly accurate screening test produces many false positives. The 2nd has a usable model, even if they consult a formula when necessary.

Strong teaching builds schemas through \textbf{clear explanations, concrete examples, contrasting cases, correct terminology, worked reasoning, and repeated application}. 1 representation is seldom enough. A diagram can reveal a relationship, a calculation can establish precision, and a counterexample can expose where a tempting analogy breaks.

\subsection*{19. Encoding Is Not the Same as Exposure}

Reading a page creates an opportunity for memory; it does not establish that the material was processed meaningfully. Encoding is supported when the learner organizes information, connects it to prior knowledge, generates examples, explains a relation, or predicts a consequence.

The danger is to turn this observation into a requirement that every minute be visibly interactive. Quiet reading, listening, and observing a skilled demonstration can be cognitively rich when the learner is constructing an internal model. Conversely, typing a response every 30 seconds can be shallow if the response is merely copying text.

The ICAP framework distinguishes passive, active, constructive, and interactive modes as a theory of engagement and proposes that more generative engagement can support deeper learning under suitable conditions. It is a useful diagnostic lens, not a universal ranking of classroom theater [R4].
